% chapters/12/evaluasi.tex
\bagian{Asesmen dan Instrumen Penilaian}
\label{sec:asesmen-unity}

\begin{asesmenbox}
\textbf{Kedudukan asesmen bab ini.} Bab ini menuntaskan \SUBCPMK{3.1}, yaitu
bab pertama pada \CPMK{3}. Ketiga indikatornya diperiksa melalui penilaian
artefak \istilah{project} (Tugas Praktikum 09) dan soal konsep yang menuntut
perbandingan dengan blok Python.

\vspace{0.4em}
Materi Bab~XII sampai Bab~XVI termasuk bahan \textbf{Proyek Akhir} dan
\textbf{Ujian Akhir Semester} pada minggu kelima belas dan keenam belas.
\end{asesmenbox}

\subsection*{Instrumen 1: Rubrik Praktikum 09}

\begin{table}[htbp]
\centering
\caption{Rubrik penilaian Praktikum 09 beserta indikator yang diukur}
\label{tab:rubrik-praktikum-09}
{\setlength{\tabcolsep}{4pt}\footnotesize%
\begin{tabularx}{\textwidth}{@{}L{3.2cm} C{1.3cm} C{2.1cm} Y@{}}
\toprule
\headrow \textbf{Kriteria} & \textbf{Bobot} & \textbf{Indikator} & \textbf{Deskripsi capaian} \\
\midrule
Penguasaan antarmuka & 30\% & \ind{a} &
Keempat panel utama beserta kegunaannya dijelaskan dengan tepat; scene dapat
disunting melalui panel \istilah{Scene} dan \istilah{Inspector}. \\
\addlinespace
Penyusunan dan penyimpanan scene & 35\% & \ind{b} &
Scene tersimpan memuat kamera, pencahayaan, dan sekurang-kurangnya tiga objek
primitif, beserta skrip C\# yang berjalan ketika scene dijalankan. \\
\addlinespace
Padanan konsep dan kerapian project & 35\% & \ind{c} &
Tabel padanan disusun sendiri dan tepat isinya; struktur folder project rapi,
dan berkas \code{.gitignore} mendaftarkan folder yang seharusnya tidak
disimpan. \\
\bottomrule
\end{tabularx}}
\end{table}

\subsection*{Instrumen 2: Soal Konsep}

Jawablah dengan singkat. Kunci jawaban beserta indikator yang diukur disajikan
pada Tabel~\ref{tab:kunci-konsep-ch12}.

\begin{enumerate}
  \item Sebutkan empat panel utama Unity beserta kegunaan masing-masing.
  \item Jelaskan mengapa sebuah berkas \istilah{scene} tidak tampak pada
        \istilah{Hierarchy}.
  \item Jelaskan gagasan penyusunan benda pada Unity, lalu bandingkan dengan
        cara menyusun benda pada Pygame.
  \item Kapan fungsi \code{Update} dan \code{FixedUpdate} dipanggil, dan apa
        perbedaan keperluan keduanya?
  \item Sebutkan tiga konsep dari blok Python yang tetap sama persis pada
        Unity, beserta padanannya.
\end{enumerate}

\begin{table}[htbp]
\centering
\caption{Kunci soal konsep Bab XII beserta indikator yang diukur}
\label{tab:kunci-konsep-ch12}
{\setlength{\tabcolsep}{4pt}\footnotesize%
\begin{tabularx}{\textwidth}{@{}C{0.8cm} C{2.1cm} Y@{}}
\toprule
\headrow \textbf{No} & \textbf{Indikator} & \textbf{Pokok jawaban} \\
\midrule
1 & \ind{a} &
\istilah{Hierarchy} memuat benda yang ada di dalam scene; \istilah{Scene}
menampilkan susunan ruangnya; \istilah{Inspector} menampilkan komponen benda
yang dipilih; \istilah{Project window} memuat seluruh berkas project. \\
\addlinespace
2 & \ind{a} &
Karena \istilah{Hierarchy} hanya memuat isi scene yang sedang dibuka,
sedangkan berkas scene adalah berkas yang berada pada \istilah{Project window}. \\
\addlinespace
3 & \ind{b} &
Benda pada Unity adalah wadah yang diisi komponen, sehingga kemampuannya
ditentukan oleh komponen yang menempel. Benda pada Pygame adalah objek dari
kelas yang ditulis sendiri, sehingga kemampuannya ditentukan oleh kelasnya. \\
\addlinespace
4 & \ind{b} &
\code{Update} dipanggil setiap bingkai untuk masukan dan hal yang mengikuti
lama bingkai; \code{FixedUpdate} dipanggil dengan selang tetap untuk
perhitungan fisika. Pemisahan itu sama dengan pembedaan langkah berubah dan
langkah tetap pada Bagian~2.2. \\
\addlinespace
5 & \ind{c} &
Antara lain: kecepatan dalam satuan per detik dikalikan selang waktu
(\code{delta_seconds} dan \code{Time.deltaTime}); arah dinormalkan sebelum
dikalikan kecepatan; dan deteksi tabrakan berbeda dari penyelesaiannya. \\
\bottomrule
\end{tabularx}}
\end{table}
